\documentclass{article}
\usepackage{lmodern}
\usepackage{amsmath}
\usepackage{listings}
\usepackage{fullpage}
\usepackage{parskip}
\usepackage{xcolor}
\lstset{
	basicstyle=\ttfamily,
	columns=fullflexible,
	frame=single,
	breaklines=true,
	postbreak=\mbox{\textcolor{red}{$\hookrightarrow$}\space},
}
\begin{document}
\title{Experiment `p\_13\_9\_recursive\_sum\_seq' Results}
\author{\today}
\date{}
\maketitle
\textbf{Experiment outcome: }SUCCESS
\\\textbf{Bad responses: }0
\\\textbf{Responses containing}~\texttt{assume}~\textbf{: }0
\\\textbf{Responses containing}~\texttt{expect}~\textbf{: }0
\\\textbf{Resolution attempts: }1
\\\textbf{Hard fails (resolution): }0
\\\textbf{Soft fails (resolution): }0
\\\textbf{Verification attempts: }1
\section*{Problem Specification}
\textbf{Problem name: }p\_13\_9\_recursive\_sum\_seq
\\\textbf{Natural language statement: }null
\\\textbf{Method signature: }p\_13\_9\_recursive\_sum\_seq(inputs: seq<int>) returns (sum: int)
\subsection*{Ensures}
\begin{itemize}
\item \texttt{sum == sum\_array(inputs)}
\end{itemize}
\subsection*{Functional Code Given}
\begin{lstlisting}
function sum_array(inputs: seq<int>): int
  ensures sum_array(inputs) == (if |inputs| == 0 then 0 else inputs[0] + sum_array(inputs[1..]))
{
  if |inputs| == 0 then 0
  else inputs[0] + sum_array(inputs[1..])
}
\end{lstlisting}
\clearpage
\section*{GenAI interactions}
Below you will find all interactions between the `user' (program) and the `assistant' (OpenAI).
\subsection*{Program $\rightarrow$ GenAI}
\begin{lstlisting}
You are given the following task to perform in Dafny:

Write a method to satisfy the below requirements.

The signature should be:

p_13_9_recursive_sum_seq(inputs: seq<int>) returns (sum: int)

The method should respect the following contract:

ensures sum == sum_array(inputs)

The contract uses the following dafny code:

function sum_array(inputs: seq<int>): int
  ensures sum_array(inputs) == (if |inputs| == 0 then 0 else inputs[0] + sum_array(inputs[1..]))
{
  if |inputs| == 0 then 0
  else inputs[0] + sum_array(inputs[1..])
}

These function(s) must not be used in your implementation of the method.

Produce and show only the Dafny body of this method, including the curly braces that surround it. Do not show the signature nor contract. You must not use 'assume' nor 'expect' anywhere in your code.
\end{lstlisting}
\subsection*{GenAI $\rightarrow$ Program}
\textbf{System fingerprint: }null
\\\textbf{ID: }chatcmpl-DnP6d9whr2ur4dHaprxhkBuFwgFfA
\\\textbf{Model: }gpt-5.5-2026-04-23
\\\textbf{Created at: }1780666951
\\\textbf{Finish reason: }stop
\\\textbf{Usage: }434 tokens in, and 543 tokens out
\begin{lstlisting}
{
  if |inputs| == 0 {
    sum := 0;
  } else {
    var tailSum := p_13_9_recursive_sum_seq(inputs[1..]);
    sum := inputs[0] + tailSum;
  }
}
\end{lstlisting}
\section*{Final Program}
\begin{lstlisting}
function sum_array(inputs: seq<int>): int
  ensures sum_array(inputs) == (if |inputs| == 0 then 0 else inputs[0] + sum_array(inputs[1..]))
{
  if |inputs| == 0 then 0
  else inputs[0] + sum_array(inputs[1..])
}

method p_13_9_recursive_sum_seq(inputs: seq<int>) returns (sum: int)
	ensures sum == sum_array(inputs)
{
  if |inputs| == 0 {
    sum := 0;
  } else {
    var tailSum := p_13_9_recursive_sum_seq(inputs[1..]);
    sum := inputs[0] + tailSum;
  }
}
\end{lstlisting}
\section*{Total Token Usage}
\textbf{Input tokens: }434
\\\textbf{Output tokens: }543
\\\textbf{Reasoning tokens: }466
\\\textbf{Sum of `total tokens': }977
\section*{Experiment Timings}
\textbf{Overall Experiment} started at 1780666949798, ended at 1780667000909, lasting 51111ms (51.11 seconds)
\\\textbf{Iteration \#1} started at 1780666949919, ended at 1780667000908, lasting 50989ms (50.99 seconds)
\end{document}
